%--------------------------------------------------------------------
\chapter{Curve su superfici}
\setlength{\parindent}{2pt}

\begin{multicols*}{2}
    \section{Piano tangente e derivata direzionale}

    \subsection{Coordinate di una curva rispetto a una parametrizzazione regolare}

    \begin{proposition} \label{prop:coordinate_curva_parametrizzazione}
        Sia $\vec{x} : U \to \RR^3$ una parametrizzazione regolare e sia $\alpha : I \to \vec{x}(U)$
        una curva. Allora $\alpha$ si scrive come:
        \[
            \alpha(t) = \vec{x}(u(t), v(t)),
        \]
        con $u(t)$, $v(t) : I \to U$ funzioni di classe $C^\infty$.
    \end{proposition}

    \begin{proof}
        Segue immediatamente dalla Proposizione \ref{prop:parametrizzazione_è_cinf}.
    \end{proof}

    \subsection{Relazione tra il piano tangente e le velocità delle curve}

    \begin{proposition}[Il piano tangente è l'insieme delle velocità delle curve sulla superficie considerata]
        Sia $\Sigma$ una superficie. Allora vale:
        \[
            T_P \Sigma = \{ \alpha'(P) \mid \alpha : I \to \Sigma \textnormal{ curva con } P \in \alpha(I) \}.
        \]
    \end{proposition}

    \begin{proof}
        Sia $P = \vec{x}(0, 0)$, dove $\vec{x} : B_\eps \to \Sigma$ è una parametrizzazione regolare di $P$.
        È sufficiente osservare che ogni vettore tangente è una combinazione lineare della forma
        $\lambda \vec{x_u} + \mu \vec{x_v}$; allora la curva
        $\alpha(t) = \vec{x}(t \lambda, t \mu)$ ha velocità $\lambda \vec{x_u} + \mu \vec{x_v}$ in
        $P$.
    \end{proof}

    \subsection{Funzioni lisce sulla superficie e derivata direzionale}

    \begin{definition}[Funzioni $C^\infty$ sulla superficie]
        Sia $\Sigma$ una superficie. Una funzione $f : \Sigma \to \RR^n$ si
        dice di classe $C^\infty$ se per ogni parametrizzazione regolare
        $\vec{x}$ di ogni punto $P \in \Sigma$, $f \circ \vec{x}$ è di
        classe $C^\infty$.
    \end{definition}

    \begin{proposition}
        Sia $\Sigma$ una superficie. Sia $P \in \Sigma$ e venga dato
        $\xi \in T_P \Sigma$. Sia data una funzione $f : \Sigma \to \RR^n$.
        Se $\alpha$, $\beta$ sono due curve su $\Sigma$
        passanti per $P$ al tempo $0$ con $\alpha'(0) = \beta'(0) = \xi$, allora:
        \[
            (f \circ \alpha)'(0) = (f \circ \beta)'(0).
        \]
    \end{proposition}

    \begin{proof}
        Sia $\vec{x}$ una parametrizzazione regolare di $P$ rispetto a $\Sigma$.
        Allora, per la Proposizione \ref{prop:coordinate_curva_parametrizzazione},
        $\alpha(t) = (u(t), v(t))$ e $\beta(t) = (p(t), q(t))$, con $u$, $v$, $p$, $q$ lisce. \smallskip

        Pertanto:
        \[
            \begin{aligned}
                (f \circ \alpha)'(0) & = \dertime{t=0} (f \circ x)(u(t), v(t))      \\[1em]
                                     & = (f \circ x)_u u'(0) + (f \circ x)_v v'(0).
            \end{aligned}
        \]

        Osserviamo che $(u'(0), v'(0)) = (p'(0), q'(0))$, dal momento che rappresentano le coordinate
        in $U$ del vettore $\xi$.
        La tesi segue allora dal fatto
        che il membro a destra diventa $(f \circ \beta)'(0)$.
    \end{proof}

    \begin{definition}[Derivata direzionale]
        Sia $\Sigma$ una superficie. Sia $P \in \Sigma$ e venga dato
        $\xi \in T_P \Sigma$. Allora, data una funzione $f : \Sigma \to \RR^n$,
        la derivata direzionale $D_\xi f(P)$ è definita come:
        \[
            D_\xi f(P) \defeq (f \circ \alpha)'(0),
        \]
        dove $\alpha$ è una qualsiasi curva su $\Sigma$
        passante per $P$ al tempo $0$ con $\alpha'(0) = \xi$.
    \end{definition}

    %\section{Operatore forma, I e II forma fondamentale}
\end{multicols*}
